% This is the ADASS_template.tex LaTeX file, 24 Oct 2025.
% It is based on the ASP general author template file, but modified to reflect the specific
% requirements of the ADASS proceedings.

% This is the aspauthor.tex LaTeX file
% Copyright 2021, Astronomical Society of the Pacific Conference Series
% Revision:  29 November 2021

% To compile, at the command line positioned at this folder, type:
% pdflatex ADASS_template
% pdflatex ADASS_template
% This will create a file called ADASS_template.pdf

\documentclass[11pt,twoside]{article}

% Do NOT use ANY packages other than asp2023. 
\usepackage{asp2023}

\aspSuppressVolSlug
\resetcounters

% References must all use BibTeX entries in a .bibfile.
% References must be cited in the text using \citet{} or \citep{}.
% Do not use \cite{}.
% See ManuscriptInstructions.pdf for more details
\bibliographystyle{asp2021}

% The ``markboth'' line sets up the running heads for the paper.
% 1 author: "Surname"
% 2 authors: "Surname1 and Surname2"
% 3 authors: "Surname1, Surname2, and Surname3"
% >3 authors: "Surname1 et al."
% Replace ``Short Title'' with the actual paper title, shortened if necessary.
% Use mixed case type for the shortened title
% Ensure shortened title does not cause an overfull hbox LaTeX error
% See ASPmanual2010.pdf 2.1.4  and ManuscriptInstructions.pdf for more details
\markboth{Stenborg}{BIAS (Bias Images As Syntax): An AI Co-designed Esolang}

\begin{document}

\title{BIAS (Bias Images As Syntax): An AI Co-designed Esolang with Python and TypeScript Interpreters}

% Note the position of the comma between the author name and the 
% affiliation number.
% Authors surnames should come after first names or initials, eg John Smith, or J. Smith.
% Author names should be separated by commas.
% The final author should be preceded by "and".
% Affiliations should not be repeated across multiple \affil commands. If several
% authors share an affiliation this should be in a single \affil which can then
% be referenced for several author names. If only one affiliation, no footnotes are needed.
% See ManuscriptInstructions.pdf and ASP's manual2010.pdf 3.1.4 for more details
\author{T.~N.~Stenborg}
\affil{The University of Sydney, Sydney, New South Wales, Australia; \email{tstenbor@acm.org}}
%\affil{$^2$Institution Name, Institution City, State/Province, Country}

% This section is for ADS Processing.  There must be one line per author. paperauthor has 9 arguments.
\paperauthor{T.~N.~Stenborg}{tstenbor@acm.org}{0000-0002-2693-9628}{The University of Sydney}{Author1 Department}{Sydney}{New South Wales}{2006}{Australia}

% There should be one \aindex line (commented out) for each author. These are used to
% build up the author index for the Proceedings. The surname must come first, followed by
% initials. Note the use of ~ before each initial to control spacing.
% The \author entries (see above) have surname last. These \aindex entries have
% surname first.
% The Aindex.py command willl create them for you after you have constructed the \author
% The first entry should be the first author, for bold-facing the author index page-reference

%\aindex{Stenborg,~T.~N.}


\begin{abstract}
BIAS (Bias Images As Syntax) is a novel visual programming language that encodes executable logic within astronomical CCD bias frames stored as FITS images. Inspired by visual and esoteric programming languages, BIAS repurposes the structure and appearance of calibration data as a programming medium, allowing valid astronomical image files to function as programs. Instructions are represented by pixel regions occupying reserved intensity bands, while execution semantics are derived from geometric relationships between opcode and operand regions. The language supports scalar and array data structures, floating-point arithmetic, flow control, and external data supplied through FITS headers and binary-table extensions.

BIAS was developed as an exploratory project investigating both unconventional programming and the practical capabilities of contemporary AI-assisted software engineering workflows. The language design, interpreter implementation, and supporting tooling were co-developed through conversational interaction with Microsoft Copilot, providing an opportunity to evaluate the effectiveness of widely accessible consumer AI systems in the creation of a novel programming language. Beyond its value as an esoteric language, BIAS may provide a useful benchmark for assessing AI reasoning on unfamiliar programming abstractions, where successful program generation requires understanding newly defined computational rules rather than relying on established programming patterns. The work demonstrates how widely-used astronomical data formats can be repurposed for creative computing while also serving as a case study in human-AI collaborative software development.
\end{abstract}

% These lines show examples of subject index entries. At this stage these have to commented
% out, and need to be on separate lines. Eventually, they will be automatically uncommented
% and used to generate entries in the Subject Index at the end of the Proceedings volume.
% Don't leave these in! - replace them with ones relevant to your paper.
%\ssindex{FOOBAR!conference!ADASS 2019}
%\ssindex{FOOBAR!organisations!ASP}

% These lines show examples of ASCL index entries. At this stage these have to commented
% out, and need to be on separate lines. Eventually, they will be automatically uncommented
% and used to generate entries in the ASCL Index at the end of the Proceedings volume.
% The ascl.py command will scan your paper on possible code names.
% Don't leave these in! - replace them with ones relevant to your paper.
%\ooindex{FOOBAR, ascl:1101.010}


\section{Introduction}
BIAS (Bias Images As Syntax) is a new programming language that encodes logic as CCD bias frames. Inspired by Piet \citep{2024MinMa..34.S139B}, BIAS was created as an exercise in experimental visual programming. It can be used for general-purpose programming, but is classified here as an esoteric programming language (esolang) due to its unconventional visual syntax. The language was developed with assistance from Microsoft Copilot.

BIAS was not intended to be a parody (e.g., \textsc{INTERCAL}), a work of software art (e.g., Piet), or a deliberate obfuscation (e.g., MALBOLGE) \citep{2005DAC..conf..144M, 2012CPCA.book..235M, 2020SIHIRbook..691D}. Any aesthetic appeal was intended to be conceptual, in the creatively novel combination of visual programming and optical astronomy calibration.

More practically, BIAS was a low-stakes project for testing the capabilities of 2026-era AI tools for software engineering. In particular, it tested the capabilites of a pre-configured, \textit{free} AI tool, Microsoft Copilot. That class of AI tool needs no investment of time or money before use, in contrast to locally-trained large language models or fee-based services.

\section{Methods}
\subsection{Program Structure}
BIAS runs on FITS files with a 16-bit ($0\textrm{--}65\,535$) intensity range. Reserved bands in the intensity range are treated as BIAS syntax. Other pixels are ignored as background data. BIAS instructions comprise pixel regions of common mean pixel intensity.

An example BIAS program is shown in Figure \ref{bias-fig-1}. BIAS programs are intended to resemble, and store information, like CCD bias frames. BIAS programs are even valid astronomical FITS files, compatible with tools like SAOImageDS9 and fv.

BIAS supports scalar and array data structures. Its instruction set facilitates flow control, and arithemtic and vector operations with native floating-point support. Floating-point support exists in some visual esolangs (e.g., Floater), but not all (e.g., Piet).

\articlefigure[width=0.6\textwidth]{P303_f1.eps}{bias-fig-1}{BIAS program that reads two floating point values from a data sidecar and sums them. The floating point addition performed is $100.0 + 17.0 = 117.0$.}

\subsubsection{Hybrid Data Sidecar}
BIAS programs support two data sidecars. FITS ASCII headers provide simple key-value data and metadata. A \textit{BIASDATA} BINTABLE extension provides bulk numerical input such as constants, arrays and tables.

\subsubsection{Opcodes and Operands}
Opcodes specify operations. Operands provide instruction parameters. BIAS programs are organised so that opcodes appear as central pixel regions surrounded by satellite operand regions.

\subsection{Program Execution}

\subsubsection{Region Detection}
BIAS scans FITS images top-to-bottom and left-to-right. Background pixels are skipped. A 4-connectivity flood fill bounds each instruction region. Region extent and centroid are calculated, and the region classified as an opcode or operand. Scanning continues until all instruction regions have been identified.

\subsubsection{Instruction Construction}
Each BIAS opcode and operand has a dedicated mean region intensity band. Opcode regions may have associated operand regions. E.g., $100 + 17$ may use an ADD opcode and two operand regions.

Associating operands to opcodes is done by adjacency. A one-pixel gap associates an operand with an opcode. To avoid ambiguity, regions may not directly abut one another.

Operands are positional parameters for opcodes. Ordering is determined geometrically. Operands are ordered by the clockwise angular displacement of their centroids, measured from the positive x-axis centred on the opcode centroid.

Opcodes and their ordered operands are then saved as BIAS Instructions.

\subsubsection{Execution Pointer}
BIAS maintains an execution pointer that traverses program pixels. Pointer state includes the current pixel, active region and traversal direction.

The entry opcode is the region whose centroid is nearest the image's top-left corner. Ties are resolved using the operand-ordering geometry.

After executing the entry instruction, the execution pointer traverses the image. Instructions execute when opcode regions are encountered. Flow may be altered by traversal-direction changes or JUMP instructions. Execution halts at the image boundary.

\section{AI Co-design}
BIAS was conceived by the author and developed with assistance from the Microsoft Copilot desktop application. Conversational engineering was used to develop architectural elements and associated source code. Microsoft Copilot is freely available to anyone with a Microsoft account.

\section{Discussion and Conclusion}
AI systems can struggle to generate working esolang programs \citep{2025Onwrdconf..213S}. Performance on BIAS programgeneration
tasks may help distinguish between AI systems capable of reasoning about novel programming scenarios and those that rely on memorised patterns or retrieval-augmented generation.

BIAS is available at \url{github.com/tstenborg/bias-images-as-syntax}.

%\acknowledgements The ASP would like to thank the dedicated researchers who are publishing with the ASP.  It will make things a lot easier if you keep this text on the same line as the \verb"\acknowledgements" command. Place it just before the bibliography. It's optional of course.


\bibliography{P303}  % For BibTex

% if we have space left, we might add a conference photograph here. Leave commented for now.
% \bookpartphoto[width=1.0\textwidth]{foobar.eps}{FooBar Photo (Photo: Any Photographer)}

\end{document}
